IoT and IIoT sites rarely want to send raw traffic to a central detector, and a detector trained only on labelled attacks cannot recognise an attack type it has never seen. Federated learning keeps traffic local but lets any site send a poisoned update. We present FL-IDS: a LightGBM classifier names known attacks, a per-site autoencoder flags traffic unlike benign traffic, federated training uses a trust filter and trimmed-mean aggregation, and verdicts become OpenFlow~1.3 rules in an emulated network. We test the whole system on all 1,985,768 unique rows of Edge-IIoTset with every connection kept on one side of every split, after correcting three defects in the released data: a placeholder-spelling label leak (reported independently by Galal), column-shifted MITM rows, and a \texttt{tshark} version difference that zeroes a feature in live extraction. In six real multi-process runs the cascade detects 99.39 to 99.81\% of known attacks (macro-recall). Packet false alarms are 2.5 to 16.3\%, but 0.34 to 1.04\% per benign connection. On attack types withheld from training the cascade detects 0.895 to 0.944 on average, against 0.723 to 0.757 for the classifier alone. With 30\% sign-flipping sites the filter keeps the autoencoder working where plain averaging destroys it, but it admits attackers in 50 to 62\% of their rounds and fails at 50\%. Replaying all 14 attack captures ends in a verified drop rule each time, with 12 named correctly. Among 15 classifiers, tree models are interchangeable on the final 130 features (macro F1 0.979 to 0.985), while boosting leads on the 91 base features.
